% ==============================================================================
% =================================== Aufgabe 3 =================================
% ==============================================================================

\chapter{Aufgabe}
\label{chap:3 Aufgabe}
\thispagestyle{fancy}
\section*{Verschaulichen Sie den Data-Mining-Prozesses anhand eines Beispiels von Internettexten.} \par Der \uline{Data-Mining-Prozess} lässt sich als automatisierte Suche nach Mustern, Regeln und Abhängigkeiten in großen Datenmengen definieren, um daraus handlungsrelevantes Wissen zu generieren. \cite{73}\cite{74} Im Kontext von \uline{Internettexten} (\textit{\acs{z. B.} aus Foren, Blogs oder sozialen Medien}) kommt dabei häufig das \uline{Text-Mining} zum Einsatz, um unstrukturierte Daten automatisiert zu erfassen und sinnvoll zuzuordnen. \cite{75} \par Ein anschaulicher Prozess am Beispiel von Internettexten umfasst folgende Phasen:
\begin{enumerate}[label=\arabic*.), leftmargin=*]
    \item \textbf{Datenerfassung (\textit{Extraktion}):} Das Unternehmen sammelt Texte aus externen Quellen wie \uline{Internetforen (\textit{\acs{z. B.} motor-talk.de})}, \uline{Blogs} oder \uline{sozialen Netzwerken}. \cite{76}\cite{77} Ein konkretes Beispiel ist das Aufspüren von Kundenfeedback nach dem Launch eines neuen Produkts. \cite{77} Da diese Daten unstrukturiert vorliegen, müssen sie zunächst für ein \uline{Customer Data Warehouse} aufbereitet werden. \cite{78}\cite{79}
    \item \textbf{Datenaufbereitung und Bereinigung:} Internettexte sind oft durch Dialekte, Slang, Sarkasmus oder Rechtschreibfehler geprägt. \cite{80} Im Rahmen der Datenveredelung müssen diese Hindernisse überwunden werden, um eine \uline{semantische Harmonisierung} zu erreichen. \cite{81} Abkürzungen müssen korrekt interpretiert werden, um Fehlinterpretationen in der späteren Analyse zu vermeiden. \cite{80}
    \item \textbf{Mustererkennung (\textit{Eigentliches Mining}):} Hier werden Algorithmen angewendet, um Bedeutungen aus den Texten abzuleiten. \cite{73}
    \begin{itemize}
        \item \textbf{Sentimentanalyse:} Ein häufiges Verfahren ist die Analyse des „Meinungsbildes“. Hierbei wird automatisch bestimmt, ob ein Text (\textit{\acs{z. B.} ein Tweet oder Forenbeitrag}) eine \uline{positive}, \uline{negative} oder \uline{neutrale Stimmung} gegenüber einem Produkt ausdrückt. Dies geschieht entweder durch Machine-Learning-Ansätze oder wörterbuchbasierte Verfahren. \cite{80}
        \item \textbf{Clustering:} Texte mit ähnlichen Inhalten oder Stimmungen können gruppiert werden, um homogene Kundensegmente zu entdecken. \cite{82}\cite{83}
    \end{itemize}
    \item \textbf{Bewertung und Wissensgenerierung:} Die gefundenen Muster werden in handlungsorientiertes Wissen transformiert. \cite{84}\cite{85}
    \begin{itemize}
        \item \textbf{Beispiel:} Erkennt das System durch Data Mining eine Häufung negativer Begriffe im Zusammenhang mit einem bestimmten Bauteil in Internetforen, kann das Unternehmen frühzeitig reagieren, bevor offizielle Reklamationen massiv zunehmen. \cite{80}
        \item \textbf{Rückgewinnungsanalyse:} Durch die Analyse der Texte können abwanderungswillige Kunden identifiziert werden, um ihnen gezielte Angebote zur Erhaltung ihrer Loyalität zu machen. \cite{86}\cite{87}
    \end{itemize}
\end{enumerate}
Zusammenfassend dient Data Mining bei Internettexten dazu, das \uline{Kundenverhalten} und \uline{die Marktposition besser zu verstehen}, indem aus einer Flut unstrukturierter Texte gezielt verwertbare Informationen extrahiert werden. \cite{88}\cite{89}

% ==============================================================================
% =================================== Aufgabe 3 =================================
% ==============================================================================